% §6 Routing evaluation (Study IV)

Study~IV performs \textbf{routing evaluation}: given the evidence in \S\ref{sec:characterization}, can pre-inference signals support appropriate model selection?
We use a calibrated logistic policy as an evaluation instrument---not a routing contribution \citep{cao2026scope}.

\subsection{Routing policy}
\label{sec:routing:policy}

Fit on CALIB: $\hat{p}(q) = P(y^{\text{opp}}{=}1 \mid c, H_w, m_w)$; tune $\tau$ for $U = \text{accuracy} - \lambda \cdot \text{cost}$; evaluate on TEST.
Assign $M_s$ if $\hat{p}(q) \geq \tau$; else $M_w$.

\subsection{Baselines}
\label{sec:routing:baselines}

Always-weak, always-strong, calibrated logistic policy, offline oracle \citep{hu2024routerbench}.
Cost: weak $=1$, strong $=3$.

\subsection{Finding 3: Routing leaves exploitable headroom}
\label{sec:routing:results}

% T4: Study IV — decision utility (fill after route-eval JSON; run per cost_lambda)
\begin{table}[t]
  \centering
  \small
  \begin{tabular}{lccccc}
    \toprule
    \textbf{Policy} & $\boldsymbol{\lambda}$ & \textbf{Acc.} & \textbf{Avg.\ cost} & $\boldsymbol{U}$ & \textbf{Opp.\ recall} \\
    \midrule
    Always-weak    & --- & TBD & TBD & TBD & TBD \\
    Always-strong  & --- & TBD & TBD & TBD & TBD \\
    Learned router & 0    & TBD & TBD & TBD & TBD \\
    Learned router & 0.05 & TBD & TBD & TBD & TBD \\
    Learned router & 0.10 & TBD & TBD & TBD & TBD \\
    Oracle (bound) & --- & TBD & TBD & --- & 100\% \\
    \bottomrule
  \end{tabular}
  \caption{Study IV (RH4): decision utility on ARC test. Router fit on CALIB; $\tau$ tuned per $\lambda$; $U=\text{acc}-\lambda\cdot\text{cost}$.}
  \label{tab:routing}
\end{table}


At $\lambda{=}0$, the calibrated policy matches always-strong (69.2\%).
The oracle reaches 74.4\%: \textbf{5.2 percentage points} of headroom, \textbf{none} exploited ($\tau{=}0$).

\paragraph{Interpretation.}
Before generation, there is useful---but incomplete---information for selecting the appropriate LLM.
Finding~2 explains part of the gap: policies that conflate \textbf{model uncertainty} with \textbf{escalation potential} cannot select appropriately.
Detection (${\sim}0.61$ AUROC) exceeds assignment under this policy class \citep{plaut2024chat,hu2024routerbench}.
